\documentclass[]{article}

\usepackage{amsmath}
\usepackage{multirow}
\usepackage{natbib}
\usepackage{graphicx} 


\begin{document}

This paper addresses the central challenge in solid sorbent design for CO₂ capture: the need to simultaneously satisfy the conflicting requirements of favorable reaction thermodynamics, high mechanical durability, and robust thermal stability. To navigate these trade-offs, we developed and applied a high-throughput computational screening methodology to systematically evaluate a vast chemical space of potential sorbent materials.

Our approach began by identifying 889 unique metal oxide-carbonate reaction pairs from the Materials Project database, filtered to include only thermodynamically stable or near-stable compounds. Each candidate was then assessed against a comprehensive set of performance metrics. The Gibbs free energy of reaction was calculated to evaluate thermodynamic viability for reversible cycling under flue gas conditions. The percentage volumetric expansion upon carbonation was quantified to predict mechanical integrity. Finally, the Tamman temperature was used as a proxy to estimate the material's resistance to thermal sintering during high-temperature operation.

The screening revealed several key findings. We found that simple binary oxides are generally poor candidates, as they occupy thermodynamic extremes. Alkali and alkaline earth oxides bind CO₂ too strongly for energy-efficient regeneration, while many transition metal oxides are unreactive. Furthermore, our analysis identified catastrophic volumetric expansion as a dominant and often overlooked failure mode. The majority of materials exhibited volume changes far exceeding 20\%, a threshold for mechanical stability, with only 14 of the 889 candidates meeting this stringent criterion. This result underscores the critical importance of dimensional stability for long-term sorbent performance.

From this work, we have learned that the materials best suited to resolve these competing demands are not simple compounds but are overwhelmingly complex, mixed-metal polyanionic frameworks. Top-performing candidates, such as sodium titanium phosphate (NaTiPO₅) and lithium vanadium phosphates, emerged from our screening by demonstrating a compelling balance of all three properties: moderate thermodynamics for reversible cycling, manageable volume expansion, and high predicted thermal stability. The isostructural substitution analysis further provided a rational design principle, showing that reaction thermodynamics can be systematically tuned by altering cation chemistry. By providing a data-driven framework for materials discovery, this study pinpoints a new and promising class of polyanionic materials for the rational design of durable, next-generation solid sorbents for industrial CO₂ capture.

\end{document}
                